% =====================================================================
%  THÈME 4 — Réactif limitant et avancement
%  Niveau DIFFICILE — Corrigé
% =====================================================================
\documentclass[11pt,a4paper]{article}
\usepackage[utf8]{inputenc}
\usepackage[T1]{fontenc}
\usepackage{lmodern}
\usepackage[french]{babel}
\usepackage{amsmath,amssymb,mathtools}
\usepackage{icomma}
\usepackage[version=4]{mhchem}
\usepackage{siunitx}
\sisetup{output-decimal-marker={,},per-mode=symbol}
\usepackage{xcolor}
\usepackage{tikz}
\usepackage[most]{tcolorbox}
\usepackage{enumitem}
\usepackage{array}
\usepackage{fancyhdr}
\usepackage[a4paper,margin=2.1cm,headheight=26pt,headsep=8pt]{geometry}

\definecolor{primary}{RGB}{142,93,168}
\definecolor{primarydark}{RGB}{82,45,108}
\definecolor{corrcol}{RGB}{176,110,20}
\newcommand{\nq}[1]{n_{\text{#1}}}
\newcommand{\Vm}{V_{\mathrm m}}

\pagestyle{fancy}\fancyhf{}
\fancyhead[L]{\small\textcolor{primarydark}{\textbf{Fiche 5} — Thème 4 : Limitant}}
\fancyhead[R]{\small\textcolor{corrcol}{\textsc{Corrigé — Difficile}}}
\fancyfoot[C]{\small\thepage}
\renewcommand{\headrulewidth}{0.8pt}
\renewcommand{\headrule}{\hbox to\headwidth{\color{primary}\leaders\hrule height \headrulewidth\hfill}}

\newtcolorbox[auto counter]{cor}[1][]{enhanced,breakable,boxrule=1pt,arc=3pt,
  left=3mm,right=3mm,top=2.4mm,bottom=2.4mm,colback=corrcol!6,colframe=corrcol,
  fonttitle=\bfseries,coltitle=white,
  title={Corrigé~\thetcbcounter\ \ifx\relax#1\relax\relax\else\textmd{--- \itshape #1}\fi}}

\newcommand{\rep}[1]{\par\smallskip\noindent\textbf{\color{corrcol}$\Rightarrow$ #1}}

\setlength{\parindent}{0pt}
\setlength{\parskip}{3pt}

\begin{document}

\begin{center}
{\LARGE\bfseries\color{primarydark}Thème 4 — Réactif limitant}\\[1pt]
{\large\color{corrcol}\textbf{Corrigé — Niveau Difficile}}\\
{\color{primary}\rule{0.6\linewidth}{1pt}}
\end{center}

\begin{cor}[lecture d'un graphe d'avancement]
\textbf{a)} À $\xi=0$ : $\nq{A,0}=\SI{0,20}{\mole}$ et $\nq{B,0}=\SI{0,15}{\mole}$
(ordonnées à l'origine des deux réactifs).\\
\textbf{b)} La courbe de \ce{A} atteint zéro en $\xi=\SI{0,10}{\mole}$, avant \ce{B} : la réaction
s'arrête là. Donc $\xi_{\max}=\SI{0,10}{\mole}$ et \boxed{\text{\ce{A} est limitant}}.\\
\textbf{c)} À $\xi_{\max}=0{,}10$ : $\nq{A}=0$ ; $\nq{B}=0{,}15-1\times0{,}10=\SI{0,05}{\mole}$
(reste, donc \ce{B} en excès) ; $\nq{C}=3\times0{,}10=\SI{0,30}{\mole}$.\\
\textbf{d)} La courbe de \ce{C} est la plus raide : sa pente $+3$ est la plus grande en valeur absolue,
car son coefficient stœchiométrique (3) est le plus grand. Les pentes valent $-2$ (\ce{A}),
$-1$ (\ce{B}) et $+3$ (\ce{C}) : \emph{la pente d'une espèce est $\pm$ son coefficient}.
\end{cor}

\begin{cor}[le cas particulier stœchiométrique]
\textbf{a)} $n(\ce{Al})=\dfrac{0{,}36}{27}=\SI{0,01333}{\mole}$ ;
$n(\ce{H2})=\dfrac{0{,}04}{2}=\SI{0,0200}{\mole}$.\\
\textbf{b)} $\dfrac{n(\ce{Al})}{2}=\dfrac{0{,}01333}{2}=0{,}006667$ ;
$\dfrac{n(\ce{H2})}{3}=\dfrac{0{,}0200}{3}=0{,}006667$. Les deux rapports sont \textbf{égaux} !\\
\textbf{c)} \textbf{Non}, il a tort : comme les rapports $n/\nu$ sont identiques, le mélange est
\emph{exactement stœchiométrique}. Aucun réactif n'est en excès et \emph{aucun} n'est « le »
limitant : ils s'épuisent \textbf{simultanément}.\\
\textbf{d)} $\xi_{\max}=0{,}006667$~mol ; $\nq{AlH3}=2\xi_{\max}=n(\ce{Al})=\SI{0,01333}{\mole}$, soit
$m=0{,}01333\times30=\boxed{\SI{0,40}{\gram}}$.
\rep{Piège classique d'annales : des rapports égaux $\Rightarrow$ pas de réactif limitant unique.}
\end{cor}

\begin{cor}[données mêlées : masse, mole et solution]
\textbf{a)} Cobalt : $n(\ce{Co})=\dfrac{58{,}93}{58{,}93}=\SI{1,00}{\mole}$.\quad
\ce{KNO3} : donné, $\SI{1,500}{\mole}$.\quad
\ce{HCl} : $n=C\times V=6\times0{,}020=\SI{0,120}{\mole}$
(\textbf{\SI{20}{\milli\litre}$=\SI{0,020}{\litre}$}).\\
\textbf{b)} $\dfrac{n(\ce{Co})}{3}=\dfrac{1}{3}\approx0{,}333$ ;\quad
$\dfrac{n(\ce{KNO3})}{2}=\dfrac{1{,}5}{2}=0{,}750$ ;\quad
$\dfrac{n(\ce{HCl})}{8}=\dfrac{0{,}120}{8}=0{,}015$.\\
\textbf{c)} Le plus petit rapport est de loin celui de \ce{HCl} ($0{,}015$) :
\boxed{\text{\ce{HCl} est le réactif limitant}}. Le cobalt et le nitrate sont en large excès.
\rep{On retrouvera ce limitant au Thème 5 pour calculer les masses et volumes de produits.}
\end{cor}

\end{document}
